


\acknowledgements{
\addtocontents{toc}{\vspace{1em}}

\begin{flushright}
\textit{Luis Fernando Delvalle}
\end{flushright}

Expreso mi agradecimiento a la Agencia de Cooperación
Internacional de Corea (KOICA) por el apoyo brindado
en el marco del proyecto de cooperación que permitió
impulsar el desarrollo de EpilexIA.

Asimismo, agradezco a la Facultad Politécnica de la
Universidad Nacional de Asunción (FP-UNA) y al Instituto
de Previsión Social (IPS) por su participación en esta
iniciativa, orientada a la aplicación de tecnologías
innovadoras en el ámbito de la salud.

Este proyecto representó una valiosa oportunidad para
fortalecer mis conocimientos y contribuir al desarrollo
de una herramienta informática destinada a apoyar
el análisis de señales electroencefalográficas.

}

\clearpage



\begingroup

\thispagestyle{empty}

\begin{center}
    {\huge\textit{Agradecimientos}\par}
\end{center}

\vspace{1em}

\begin{flushright}
\textit{Rodrigo Amarilla Sanabria}
\end{flushright}


En primer lugar, agradezco a Dios por darme la fortaleza,
la perseverancia y la sabiduría necesarias para afrontar
los desafíos que se presentaron a lo largo de mi formación
académica y durante el desarrollo de este trabajo.

A mis padres y hermanos, por su cariño, apoyo incondicional
y por acompañarme en las distintas etapas de mi vida.
A mi amada esposa, Paola Páez, por su amor, comprensión, paciencia
y por estar a mi lado durante este camino, especialmente
en los momentos de mayor dificultad.

Expreso mi sincero agradecimiento al profesor Diego Pinto,
tutor de esta tesis, por su orientación, confianza,
dedicación y acompañamiento durante el desarrollo del
proyecto EpilexIA. Sus conocimientos y recomendaciones
fueron fundamentales para la realización de este trabajo.


A mi compañero de tesis, Luis Fernando Delvalle,
por el esfuerzo compartido, el compromiso y el
trabajo en equipo durante el desarrollo de EpilexIA y,
sobre todo, por su calidad humana, por
escucharme con paciencia, brindarme sus consejos
y darme palabras de aliento durante este camino.
Su amistad, comprensión y apoyo fueron muy
valiosos para encontrar la fortaleza necesaria
para seguir adelante y no rendirme ante las
adversidades. Asimismo,
agradezco a Iván Ríos, Ignacio Lezcano, Edgar Ayala,
Luis Cubilla, Eladio Quintana y Maximiliano Villamayor,
integrantes del equipo de EpilexIA, por su colaboración,
conocimientos y aportes al desarrollo de esta iniciativa.

A la Agencia de Cooperación Internacional de Corea (KOICA),
por el apoyo brindado al proyecto de cooperación en cuyo
marco se desarrolló EpilexIA. De igual manera, agradezco
a la Facultad Politécnica de la Universidad Nacional de
Asunción (FP-UNA) y al Instituto de Previsión Social (IPS)
por su participación y contribución a esta iniciativa de
innovación tecnológica aplicada al ámbito de la salud.


\needspace{6\baselineskip}
A mi amigo Fernando Moreno, por haberme tendido
la mano en uno de los momentos más difíciles
de mi vida, cuando más necesitaba apoyo y una
oportunidad para salir adelante.
Su generosidad y confianza en mí fueron
fundamentales para superar aquella etapa
y continuar persiguiendo mis objetivos.
Siempre llevaré conmigo una profunda gratitud
por aquel gesto de amistad y solidaridad.


Finalmente, expreso mi gratitud a todas las personas que,
de una u otra manera, contribuyeron a mi formación y a
la culminación de este trabajo, que representa una etapa
muy importante de mi vida personal y profesional.


\par
\endgroup


